% Pattern: proof/amsthm/amsart. A remark between a statement and its proof,
% as \end{lemma} \begin{remark} We will later apply ... \end{remark}
% \begin{proof}, does not take the proof, unless the statement cites another
% paper; a proof inside a remark proves the remark, and so does each proof
% after it there; and a display in the text between a statement and its proof
% does not break them.
% seen in the gold pilot (set X)
% want lem-1 statement/lemma 1
% want rem-1 remark 2
% want proof-1 proof - proves=lem-1 flag=remark-between
% want lem-2 statement/lemma 3
% want rem-2 remark 4
% want proof-2 proof - proves=rem-2
% want rem-3 remark 5
% want proof-3 proof - proves=rem-3
% want proof-4 proof - proves=rem-3
% want lem-3 statement/lemma 6
% want proof-5 proof - proves=lem-3 flag=text-between
\documentclass{amsart}
\newtheorem{lem}{Lemma}
\theoremstyle{remark}
\newtheorem{rem}[lem]{Remark}
\begin{document}
\begin{lem} Each tree with two vertices or more has a leaf. \end{lem}

\begin{rem} We will later apply this to the trees of a forest. \end{rem}

\begin{proof} Take a longest path. \end{proof}

\begin{lem}\cite{X} Each finite graph has a spanning forest. \end{lem}
\begin{rem} Each forest with $k$ trees and $n$ vertices has $n-k$ edges. \end{rem}
\begin{proof} Count the edges of each tree. \end{proof}

\begin{rem} We list two facts.
\begin{proof} The first is clear. \end{proof}
\begin{proof} The second is clear too. \end{proof}
\end{rem}

\begin{lem} Each tree has a proper two-colouring. \end{lem}

We note a formula before the proof.
\begin{equation}
 \chi(T) \le 2.
\end{equation}
It is the bound.

\begin{proof} Colour by the parity of the distance to a root. \end{proof}
\end{document}
